\section{Clock dilution and the whole-operator obstruction}
\label{sec:dilution}

Let $\Hcal_{\mathrm{acc}}$ be the closed span of vertices on the prime cycles.
It is a reducing subspace, and
\[
L_s\big|_{\Hcal_{\mathrm{acc}}}=\bigoplus_p B_{p,s}.
\]
The next theorem uses only disjointness, nonnegative clocks, the exact total,
and
$\log p/\ell(p)\to0$.  It is therefore independent of all modeling choices
about how the source clock is distributed.

\begin{theorem}[Clock-dilution obstruction]
\label{thm:noncompact}
For every $s\in\C$ with $\sigma=\RePart s>0$ and every nonnegative exact-clock
allocation satisfying \eqref{eq:exact-clock}, the accepted restriction of
$L_s$ is noncompact and
\begin{equation}
\norm{L_s}_{\ess}=1.
\label{eq:essential-norm}
\end{equation}
Consequently the whole source-weighted vertex adjacency is noncompact.
\end{theorem}

\begin{proof}
Choose on every $\Gamma_p$ an edge $e_p$ of minimum roof.  Then
\[
\tau(e_p)\le\frac{\log p}{\ell(p)}.
\]
Let $u_p$ be its source.  Distinct prime cycles are vertex-disjoint, so
$(\delta_{u_p})_p$ is an orthonormal, hence weakly null, sequence.  Yet
\begin{equation}
\norm{L_s\delta_{u_p}}
=e^{-\sigma\tau(e_p)}
\ge p^{-\sigma/\ell(p)}\longrightarrow1.
\label{eq:edge-witness}
\end{equation}
A compact operator maps weakly null bounded sequences to norm-null
sequences, proving noncompactness.

All nonnegative-roof edge weights have modulus at most one, and the functional
graph has at most one outgoing edge per vertex, so $\norm{L_s}\le1$.  After
subtracting any compact operator,
\eqref{eq:edge-witness} still supplies a tail with lower norm tending to one.
Thus $\norm{L_s}_{\ess}\ge1$, while the operator-norm upper bound gives
\eqref{eq:essential-norm}.
\end{proof}

\begin{theorem}[No finite Schatten repair]
\label{thm:no-schatten}
For every $q>0$ and $\sigma>0$, the accepted restriction of $L_s$ does not
belong to $\Sq$.
\end{theorem}

\begin{proof}
By
\eqref{eq:block-singular}, its $q$th singular-value sum on the $p$-block is
\[
\sum_{e\in\Gamma_p}e^{-q\sigma\tau(e)}.
\]
The convexity of $x\mapsto e^{-q\sigma x}$ and the exact total give
\begin{equation}
\sum_{e\in\Gamma_p}e^{-q\sigma\tau(e)}
\ge \ell(p)\exp\!\left(-\frac{q\sigma\log p}{\ell(p)}\right).
\label{eq:jensen-schatten}
\end{equation}
The right side tends to infinity.  In particular the block contributions do
not tend to zero, so their sum over primes diverges.  This is precisely the
$q$th power sum of the singular values of the orthogonal direct sum.
\end{proof}

Trace class is the $q=1$ case.  Therefore the ordinary Hilbert-space Fredholm
determinant $\det_{\Hcal}(I-zL_s)$ requested by the same-object protocol is
not defined on any half-plane through the usual trace-class theory
\citep{Simon1977}.  The failure is stronger than divergence of a particular
rank-one edge expansion: the operator is not even compact.

\begin{theorem}[Essential approximate unit circle]
\label{thm:essential-circle}
For every $s$ with $\RePart s>0$,
\begin{equation}
\T\subset\sigma_{\ap,\ess}(L_s).
\label{eq:essential-circle}
\end{equation}
Hence $I-zL_s$ is not Fredholm whenever $|z|=1$.
\end{theorem}

\begin{proof}
Fix $\lambda\in\T$.  By
\cref{lem:cycle-block}, the eigenvalues of the prime block are the
$\ell(p)$ roots of $p^{-s}$.  Their common modulus tends to one by
\eqref{eq:block-radius}; their angular mesh is $2\pi/\ell(p)\to0$, and the
common phase offset has size $O(\log p/\ell(p))\to0$.  Choose an exact
block eigenvalue $\lambda_p$ tending to $\lambda$ and a normalized
eigenvector $x_p$ supported on that block.  The vectors lie in mutually
orthogonal finite-dimensional blocks, hence $x_p\rightharpoonup0$, while
\[
\norm{(L_s-\lambda I)x_p}=|\lambda_p-\lambda|\longrightarrow0.
\]
This singular Weyl sequence proves
\eqref{eq:essential-circle}.  If $|z|=1$, take $\lambda=z^{-1}$; a Fredholm
operator cannot admit such a singular sequence for its zero spectral value.
\end{proof}

The theorem does not identify a self-adjoint spectrum and does not concern
Riemann zeros.  It says that the graph-step verifier creates essential
spectrum at the unit circle before any critical-strip spectral mechanism is
available.

\subsection{An exact criterion behind the example}

\begin{proposition}[Disjoint cycles with prescribed totals]
\label{prop:compactness-criterion}
Let $B=\bigoplus_a B_a$ be a weighted adjacency on pairwise disjoint finite
cycles of lengths $\ell_a$, with nonnegative edge roofs $\tau_{a,j}$ and
totals $T_a=\sum_j\tau_{a,j}$.  For a fixed allocation, $B$ is compact if
and only if
\[
\min_j\tau_{a,j}\longrightarrow\infty.
\]
There exists some allocation with the prescribed totals for which $B$ is
compact if and only if
\begin{equation}
\frac{T_a}{\ell_a}\longrightarrow\infty.
\label{eq:allocation-criterion}
\end{equation}
\end{proposition}

\begin{proof}
The singular values of each block are its edge-weight moduli.  A block direct
sum of finite matrices is compact exactly when its largest block singular
value tends to zero.  Since that value is
$e^{-\sigma\min_j\tau_{a,j}}$, this gives the first equivalence.  Necessity
of
\eqref{eq:allocation-criterion} follows from
$\min_j\tau_{a,j}\le T_a/\ell_a$.  For sufficiency, use the uniform
allocation $\tau_{a,j}=T_a/\ell_a$.
\end{proof}

The prime verifier has
$T_p/\ell(p)=\log p/\ell(p)\to0$, the opposite of the compactness regime.
This criterion is included to expose the mechanism; we make no claim that
its abstract weighted-shift content is novel.
